% === Begin Label Data ===
% ID: 371
% 难度星级: 
% 标签: 
% 备注: 
% 组卷引用次数: 0
% === End  Label Data ===

\begin{problem}{2016}{G}{四川卷（理）}{15}{解析几何}
在平面直角坐标系中，当 $ P(x,y) $ 不是原点时，定义 $ P $ 的“伴随点”为 $ P'\left(\dfrac{y}{x^2+y^2},\dfrac{-x}{x^2+y^2}\right) $；当 $ P $ 是原点时，定义 $ P $ 的“伴随点”为它自身。平面曲线 $ C $ 上所有点的“伴随点”所构成的曲线 $ C' $ 定义为曲线 $ C $ 的“伴随曲线”，现有下列命题：

\circled{1} 若点 $ A $ 的“伴随点”是点 $ A' $，则点 $ A' $ 的“伴随点”是点 $ A $；

\circled{2} 单位圆的“伴随曲线”是它自身；

\circled{3} 若曲线 $C$ 关于 $x$ 轴对称，则其“伴随曲线” $C'$ 关于  $y $ 轴对称;

\circled{4} 一条直线的“伴随曲线”是一条直线.

其中的真命题是 \underline{\hspace{4em}}.（写出所有真命题的序列）
\end{problem}

\begin{answer}

\end{answer}

\begin{solutions}

\end{solutions}